% P03.1 rational-defect terminal boundary core proof, live integration.

\begin{proof}[\Cref{thm:p031-rational-defect-terminal-boundary-core}]
Write $h:=g_v(x)$.  If $h\leq0$, the set in
\eqref{eq:p031-rational-defect-rooted-terminal-trace} is empty for every
$\xi$, so $\rho_{W,v,x}=0$ and
\eqref{eq:p031-rational-defect-rooted-terminal-maximum} is immediate.
Suppose henceforth that $h>0$.

For $0<\alpha<h$, the set
$B_{\alpha,W,v}(x)$ is nonempty by
\Cref{cor:p031-rational-defect-coherent-carrier}.  It is compact because it
is closed in the compact set $\partial W$.  If
$0<\alpha<\beta<h$, then
\[
 \Gamma_{\beta,W}(x)\subset\Gamma_{\alpha,W}(x).
\]
Indeed, the set on the left is a connected subset of
$\{g_v\geq\alpha\}$ that contains $x$, so it lies in the component of $x$
there.  Therefore
\begin{equation}
\label{eq:p031-rational-defect-rooted-contact-nesting}
             B_{\beta,W,v}(x)\subset B_{\alpha,W,v}(x).
\end{equation}
The nested nonempty compact subsets $B_{\alpha,W,v}(x)$ of the compact
boundary have nonempty intersection.  Hence
$\mathfrak T_{W,v}(x)$ is nonempty and compact.

The rooted trace is bounded above by $h$.  For $0<r\leq h$, nesting gives
\begin{equation}
\label{eq:p031-rational-defect-rooted-trace-superlevel}
 \{\xi\in\partial W:\rho_{W,v,x}(\xi)\geq r\}
   =\bigcap_{0<\alpha<r}B_{\alpha,W,v}(x),
\end{equation}
while the left-hand side is empty for $r>h$.  Thus every positive closed
superlevel of $\rho_{W,v,x}$ is closed, and the function is upper
semicontinuous.  Every point of $\mathfrak T_{W,v}(x)$ has rooted trace
$h$, whereas no rooted trace exceeds $h$.  This proves
\eqref{eq:p031-rational-defect-rooted-terminal-maximum} and identifies its
maximizer set as in
\eqref{eq:p031-rational-defect-terminal-argmax}.

If $\xi\in B_{\alpha,W,v}(x)$, every ordinary neighborhood of $\xi$ meets
$\Gamma_{\alpha,W}(x)$, on which $(g_v)_+\geq\alpha$.  Hence
$\tau_{W,v}(\xi)\geq\alpha$.  Taking the supremum over all levels selected
by $\xi$ proves
\eqref{eq:p031-rational-defect-rooted-trace-bound}.  In particular, every
$\xi\in\mathfrak T_{W,v}(x)$ satisfies
$\tau_{W,v}(\xi)\geq h$, completing
\eqref{eq:p031-rational-defect-terminal-argmax}.

Finally, if $h=\sup_W(g_v)_+$, then
\Cref{thm:p031-rational-defect-boundary-trace} gives
$\tau_{W,v}\leq h$ on $\partial W$.  Equality holds at every point of the
terminal core, proving
\eqref{eq:p031-rational-defect-terminal-max-trace}.
\end{proof}

\begin{proof}[\Cref{cor:p031-rational-defect-rooted-terminal-formula}]
For each admissible pair $(W,v)$ through $x$,
\Cref{thm:p031-rational-defect-terminal-boundary-core} gives
\[
       (v(x)-s_{\mathbf Q,X}(x))_+
          =\max_{\xi\in\partial W}\rho_{W,v,x}(\xi).
\]
Take the supremum over the same relatively compact test domains and local
rational competitors as in
\eqref{eq:p031-rational-defect-end-height-formula}.
\end{proof}

\begin{proof}[\Cref{cor:p031-finite-flat-persistent-terminal-anchor}]
Each $K_\alpha$ is compact and connected: it is the closure of the connected
set $\Gamma_{\alpha,W}(x)$ inside the compact set $\overline W$.  It contains
$x$ and $\xi$.  The nesting argument in
\eqref{eq:p031-rational-defect-rooted-contact-nesting} also gives
\begin{equation}
\label{eq:p031-rational-defect-rooted-carrier-nesting}
                    K_\beta\subset K_\alpha
                    \qquad(0<\alpha<\beta<h).
\end{equation}

For each $\alpha$, apply
\Cref{lem:p031-proper-open-compact-component-anchor} to
$\pi^{-1}(K_\alpha)\to K_\alpha$, with anchor point $x$.  The component
$\widetilde K_\alpha(\eta)$ is compact and connected and meets the finite
fiber $\pi^{-1}(x)$.  By
\eqref{eq:p031-rational-defect-rooted-carrier-nesting},
$\widetilde K_\beta(\eta)$ is a connected subset of
$\pi^{-1}(K_\alpha)$ containing $\eta$, and therefore
\[
             \widetilde K_\beta(\eta)
                  \subset\widetilde K_\alpha(\eta)
             \qquad(0<\alpha<\beta<h).
\]
Consequently the sets
\[
 A_\alpha:=\widetilde K_\alpha(\eta)\cap\pi^{-1}(x)
\]
form a decreasing family of nonempty subsets of the finite set
$\pi^{-1}(x)$.  Their intersection is nonempty, which is exactly
\eqref{eq:p031-finite-flat-persistent-terminal-anchor}.  Any point $e$ in
that intersection is the asserted common anchor.

The trace identity
\eqref{eq:p031-finite-flat-boundary-trace-pullback} and
\eqref{eq:p031-rational-defect-terminal-argmax} give
\eqref{eq:p031-finite-flat-persistent-terminal-trace}.
\end{proof}
